\documentclass[11pt,a4paper]{article}
% geometry.sty is unavailable in this MiKTeX install (no package server);
% equivalent 2.3cm margins set with plain LaTeX lengths instead.
\setlength{\oddsidemargin}{-0.24cm}
\setlength{\evensidemargin}{-0.24cm}
\setlength{\textwidth}{16.4cm}
\setlength{\topmargin}{-0.24cm}
\setlength{\headheight}{0cm}
\setlength{\headsep}{0cm}
\setlength{\textheight}{25.1cm}

% Section + page numbers are a stated marking requirement (guidelines §4/§6).
\pagestyle{plain}

\title{Visual Analytics of the 1994 UCI Adult Census Dataset:\\
Bayesian Imputation, Dimensionality Reduction and an Interactive Tableau Dashboard}
\author{Visual Analytics --- Resit Coursework}
\date{July 2026}

\begin{document}
\maketitle

\begin{abstract}
\noindent We present an exploratory visual analysis of the 1994 UCI Adult
census extract (48{,}842 records, 14 attributes) for a policy-analyst user,
combining Bayesian multinomial imputation of missing employment attributes
with two dimensionality-reduction projections (PCA and t-SNE) delivered as
an interactive Tableau dashboard. The main finding is that reported
\emph{capital income}, not wage-related attributes, is the sharpest visual
separator of high earners: the PCA projection splits the population into
three strata by capital-gain/loss activity, and within the capital-gain
stratum 61.4\% of respondents earn over \$50K versus 18.9\% in the
no-capital-activity majority (base rate 23.9\%). Within every stratum,
income rises along a work/status composite (hours worked, education, age,
sex). Occupation and education gradients are steep --- from 3.7\%
(\textit{Other-service}) to 47.6\% (\textit{Exec-managerial}) and from
15.9\% (high-school) to 74.0\% (professional school) --- and the male
high-income rate (30.4\%) is nearly triple the female rate (10.9\%).
\end{abstract}

\section{Introduction}\label{sec:intro}
The Adult dataset, extracted from the 1994 US Census by Kohavi~\cite{kohavi},
records demographic and employment attributes together with a binarised
income label ($\leq$\$50K / $>$\$50K). It is a standard benchmark for
classification, but classification accuracy alone tells an analyst little
about \emph{structure}: which population segments drive high income, how
those segments overlap, and where the data itself is unreliable.

The intended user of this work is a \textbf{policy analyst} studying the
socio-economic drivers of income inequality: a domain expert who is fluent in
census vocabulary and comfortable with interactive dashboards, but who does
not read code and should not have to trust opaque model outputs. Following
Munzner's nested model~\cite[Ch.~4]{munzner}, this domain situation drives
every downstream choice: the task abstraction (\S\ref{sec:tasks}), the visual
encoding idioms (\S\ref{sec:vis}), and the algorithms (PCA, t-SNE, NUTS).

Why visualization rather than summary tables? Anscombe's Quartet
\cite[pp.~8--9]{munzner} is the canonical argument: datasets with identical
summary statistics can differ radically in structure that is obvious the
moment it is plotted. The census equivalent: a table of income-by-education
means cannot reveal that capital-gain reporters form a separate,
majority-high-income stratum of the population --- the PCA projection makes
this visible in a single view (\S\ref{sec:vis}).

The deliverable is a packaged Tableau workbook
(\texttt{adult\_dashboard.twbx}) with two projection views
(\texttt{Proj\_PCA}, \texttt{Proj\_tSNE}), three group-comparison views with
an absolute/proportion toggle (\texttt{Income\_by\_Education},
\texttt{Income\_by\_Occupation}, \texttt{Income\_by\_Age}), an imputation
quality view (\texttt{Bayesian\_Imputation\_QA}), and an overview dashboard.

\section{Data Preparation and Abstraction}\label{sec:data}
The analysis dataset combines the UCI \texttt{adult.data} and
\texttt{adult.test} files (48{,}842 rows, 14 features, one label).
Combining them required stripping trailing periods from test-set labels
and leading whitespace from all categorical values. The full pipeline is
reproducible: staged Python scripts with printed verification gates, fixed
random seeds, and all dimensionality-reduction hyperparameters logged to
\texttt{params.json}.

\paragraph{Anomalies, quantified.}
The \texttt{?} sentinel marks missing values in exactly three columns:
\texttt{workclass} (2{,}799), \texttt{occupation} (2{,}809) and
\texttt{native-country} (857). The missingness is not homogeneous, and we
name the two kinds separately. All 2{,}799 \texttt{workclass} gaps co-occur
with \texttt{occupation} gaps (structural co-missingness of the employment
block), and the 10 remaining \texttt{occupation} gaps belong to
\texttt{Never-worked} respondents for whom no occupation exists ---
these are set directly to an explicit \emph{None (never worked)} category
rather than being modelled. Only \texttt{native-country} exhibits ordinary
item missingness; its 857 gaps are recoded to an explicit \emph{Unknown}
category because a 41-class Bayesian model is disproportionate for a column
that is not a projection input --- a stated trade-off, not a hidden one.
There are 52 full-row duplicates, retained as plausible identical census
records. \texttt{education} and \texttt{education-num} are a redundant 1:1
pair; the ordinal \texttt{education-num} is kept for modelling and the text
form retained only as tooltip context (a \emph{derive} decision in Munzner's
sense~\cite[pp.~49--53]{munzner}).

\paragraph{Data abstraction.}
In Munzner's Ch.~2 vocabulary the dataset is a flat table of 48{,}842 items.
Quantitative attributes: \texttt{age}, \texttt{fnlwgt},
\texttt{capital-gain}, \texttt{capital-loss}, \texttt{hours-per-week};
ordinal: \texttt{education-num}; categorical: \texttt{workclass},
\texttt{occupation}, \texttt{marital-status}, \texttt{relationship},
\texttt{race}, \texttt{sex}, \texttt{native-country}; the label
\texttt{income} is \emph{binary} categorical, which matters for the colour
encoding (\S\ref{sec:vis}). Categorical/ordinal attributes act as keys and
quantitative attributes as values --- exactly the Dimensions/Measures split
Tableau inherits from Polaris~\cite[p.~40]{munzner}.

\paragraph{Bayesian imputation.}
The employment attributes were imputed with sequential Bayesian multinomial
logistic regressions in PyMC (NUTS sampler via nutpie; 1{,}000 posterior
draws, 500 tuning steps, 4 chains, fixed seed): first \texttt{workclass},
then \texttt{occupation} with the (now complete) \texttt{workclass} as an
additional predictor, mirroring their structural dependency. Predictors ---
standardised \texttt{age}, \texttt{education-num},
\texttt{hours-per-week}, \texttt{sex}, and marital status collapsed to
married / previously married / never married --- were chosen as complete
columns plausibly associated with employment type. Each model was fitted on
a stratified subsample of $\approx$8{,}000 observed rows (minimum 15 rows
per class) to keep full NUTS sampling tractable; predictions cover every
missing row. Convergence was checked before proceeding: maximum
$\hat{R}=1.000$ (min bulk ESS 1{,}094) for the workclass model and
$\hat{R}=1.010$ (min bulk ESS 777) for the occupation model, both within
the pre-registered $\hat{R} \leq 1.05$ gate.

The standard ``posterior mean / posterior std'' imputation recipe assumes a
continuous target; for these categorical targets we store the
posterior-predictive \emph{modal category} as the imputed value, and the
modal probability plus predictive entropy as per-row uncertainty columns.
This makes the imputation's confidence itself visualisable: mean modal
probability was 0.756 for workclass but only 0.297 for occupation (14
classes), i.e.\ the model is honest that occupation is hard to recover.
A known artefact of modal point imputation is surfaced rather than hidden:
it over-assigns majority classes (99.8\% of imputed workclass rows became
\textit{Private} versus 73.6\% in the observed marginal). The
\texttt{Bayesian\_Imputation\_QA} sheet plots imputed against observed
distributions precisely so the analyst sees this shrinkage; the stored
entropy columns quantify it per row. A mode-imputation baseline (which
would place \emph{100\%} of rows in the majority class with no uncertainty
estimate) is retained in the dataset as a clearly labelled fallback
comparison only.

\paragraph{Projection feature space.}
Both projections use the same six standardised features ---
\texttt{age}, \texttt{education-num}, \texttt{hours-per-week},
$\log(1{+}\texttt{capital-gain})$, $\log(1{+}\texttt{capital-loss})$,
\texttt{sex} --- recorded verbatim in a \texttt{variables\_used} column
that every projection tooltip displays. Exclusions are deliberate:
\texttt{income} is the conjectured categorisation the projections are meant
to \emph{verify}~\cite[\S13.4.3]{munzner} and would leak the answer into
the layout; \texttt{fnlwgt} is a survey design weight describing the
sampling frame, not the respondent; the high-cardinality nominals are
surfaced through tooltips instead of being one-hot inflated into the
distance metric.

\section{Task Definition}\label{sec:tasks}
Tasks are phrased in Munzner's action/target vocabulary~\cite[Ch.~3]{munzner}
so that each view can be checked against the question it exists to answer.

\begin{enumerate}
  \item \textbf{Discover and verify cluster structure.} \emph{Discover}
  whether the six-attribute feature space contains clusters or strata, and
  \emph{verify} whether they align with the income category
  (targets: clusters, similarity). Answered by \texttt{Proj\_PCA} and
  \texttt{Proj\_tSNE}: the analyst inspects colour-coded structure, then
  uses tooltips to \emph{identify} what individual points are.
  \item \textbf{Compare income distributions across groups.} \emph{Compare}
  the income distribution (target: distribution) across education levels,
  occupations and 5-year age bands --- in absolute counts when the analyst
  asks ``where are the people?'' and proportions when asking ``which group
  has the highest rate?''. Answered by the three \texttt{Income\_by\_*}
  sheets with the \textsf{Show As} parameter toggle.
  \item \textbf{Summarise imputation quality.} \emph{Summarise} the derived
  (imputed) data and \emph{compare} imputed against observed category
  distributions, so the analyst knows how much trust to place in the
  5.7\% of rows whose employment attributes are model output. Answered by
  \texttt{Bayesian\_Imputation\_QA}.
\end{enumerate}

\section{Visualisation Justification}\label{sec:vis}
We justify each view at the idiom level of the nested
model~\cite[Ch.~4]{munzner}, citing the perceptual evidence for every
channel choice.

\paragraph{Projections as scatterplots.}
Dimensionality reduction is attribute aggregation in Munzner's
taxonomy~\cite[\S13.4.3]{munzner}, and its display follows the
document-collection pattern given there: high-dimensional table
$\rightarrow$ 2D layout $\rightarrow$ scatterplot colour-coded by a
conjectured categorisation, with the task of verifying whether visual
clusters match it. Two-dimensional scatterplots are empirically the safest
idiom for inspecting reduced data~\cite{sedlmair}; a 3D scatterplot was
rejected on both this DR-specific evidence and the general argument against
unjustified 3D for abstract data (depth is read far less accurately than
planar position; occlusion and perspective distortion break the position
and size channels)~\cite[Ch.~6.3]{munzner}.

\begin{table}[h]\centering\small
\begin{tabular}{ll}
\hline
Idiom & \texttt{Proj\_PCA} / \texttt{Proj\_tSNE} \\
\hline
What: data & Table, 48{,}842 items $\times$ 6 quantitative attributes \\
What: derived & 2 synthetic attributes (PC1/PC2 or TSNE1/TSNE2); \\
 & \texttt{variables\_used} records the contributing attributes \\
How: encode & Point marks; aligned position for the 2 synthetic attributes; \\
 & hue for binary \texttt{income}; 30\% opacity against overplotting \\
How: facet & Juxtaposed in the overview dashboard; same data, same colour \\
 & encoding, independent navigation \\
How: reduce & None (all items shown; aggregation happens in attribute space) \\
Why: task & Discover/verify income-aligned cluster structure; identify \\
 & points via details-on-demand tooltips \\
Scale & 48{,}842 marks \\
\hline
\end{tabular}
\caption{Idiom analysis of the projection sheets, following the summary-table
format Munzner uses for the equivalent DR scenario \cite[p.~319]{munzner}.}
\end{table}

Position on a common scale is the most accurate magnitude channel, so it is
reserved for the projection coordinates themselves; \texttt{income}, a
binary categorical attribute, takes hue --- the highest-ranked identity
channel after spatial position~\cite[Fig.~5.6]{munzner}. Encoding a
categorical attribute with an ordered channel (or vice versa) would violate
the expressiveness principle~\cite[p.~100]{munzner}. With 48{,}842 marks the
scatterplot is beyond the ``dozens to hundreds'' comfort zone Munzner notes
for scatterplots~\cite[\S7.4]{munzner}; 30\% mark opacity turns overlap
into readable density, which is what rescues the idiom at this scale.

Both projection layouts carry two caveats that the dashboard states and the
analyst must internalise: DR layouts are affine invariant, so only
\emph{relative} distances carry meaning, never absolute position or
orientation; and fine-grained distances are less reliable than large-scale
cluster structure because information is necessarily lost in the
reduction~\cite[pp.~319--320]{munzner}. The t-SNE view additionally shows
thread-like micro-clusters: census attributes are quasi-discrete, so many
records are exact duplicates in feature space --- an artefact of the data's
granularity, not evidence of fine social structure. t-SNE hyperparameters
(perplexity 30, 1{,}000 iterations, PCA initialisation, fixed
seed~\cite{vandermaaten}) are logged in \texttt{params.json}.

\paragraph{Tooltips.}
Every projection mark carries a details-on-demand tooltip (age, education,
occupation, income, hours/week, sex, and the \texttt{variables\_used}
string), implementing the overview-first, details-on-demand
mantra~\cite{shneiderman}\cite[\S6.7]{munzner}. Displaying
\texttt{variables\_used} in the tooltip operationalises the requirement that
DR plots communicate which original variables produced them --- the
synthetic axes are otherwise uninterpretable to the analyst.

\paragraph{Group comparison bars.}
The \texttt{Income\_by\_*} sheets are stacked bar charts: one categorical
key (region: separated and ordered along an axis) and one quantitative
value (aligned position), with hue for the income
split~\cite[\S7.5]{munzner}. Ordering is a deliberate, named decision in
each case: education is ordered by \texttt{education-num} (its inherent
ordinal order), occupations are value-sorted to reveal the gradient, and
age bands follow their natural sequence. The mandatory
absolute\,/\,proportion toggle is one \textsf{Show As} parameter driving a
calculated field (not two parallel sheets), and its justification is
perceptual: absolute counts and proportions support different comparison
tasks when group sizes differ~\cite[\S5.6]{munzner} --- \textit{Doctorate}
holders are a tiny bar in absolute view yet the second-highest rate
(72.6\%) in proportional view.

\paragraph{Colour.}
The income palette (blue $\leq$\$50K, orange $>$\$50K) was validated
computationally, not by eye: simulated protanopia/deuteranopia/tritanopia
colour distances all exceed $\Delta E = 96$ and both colours clear 3:1
contrast on the workbook background. This matters because red--green
colour-vision deficiency affects roughly 8\% of males~\cite[p.~235]{munzner},
and the two hues also differ in luminance, satisfying the ``get it right in
black and white'' rule~\cite[p.~140]{munzner}. The palette ships as
\texttt{Preferences.tps}. The QA sheet uses neutral grey for observed rows
so the imputed (blue) distribution is the visually dominant comparison.

\paragraph{Imputation QA view.}
\texttt{Bayesian\_Imputation\_QA} juxtaposes the workclass distribution of
imputed rows against observed rows as aligned bars sharing one axis ---
``eyes beat memory'': simultaneous views beat remembered
comparison~\cite[\S6.5]{munzner}. This view exists because imputation is a
\emph{derive} operation whose output the analyst must be able to audit; it
is where the modal-imputation majority-shrinkage (\S\ref{sec:data}) is
visible instead of buried in a methods paragraph.

\paragraph{Dashboard composition.}
The overview dashboard is a juxtaposed, coordinated multi-view
design~\cite[Ch.~12]{munzner}: all views share the same underlying data and
the same income colour encoding (shared encoding, all-data), while
navigation is independent per view. The parameter control is placed once,
above the views it drives. No 3D, gradients, shadows or decorative effects
are used --- function first, form next~\cite[\S6.10]{munzner}.

\section{Conclusion}\label{sec:conclusion}
\paragraph{What the analyst learns.}
The projections answer Task~1 affirmatively and specifically: the feature
space is stratified, and the strata are economic. PCA's second component is
dominated by capital income ($+0.70$ capital-loss, $-0.63$ capital-gain
loadings), splitting the population into a capital-gain stratum (8.3\% of
respondents, 61.4\% of whom earn $>$\$50K), a capital-loss stratum (50.1\%
high earners) and the no-capital-activity majority (18.9\%). Reporting
\emph{any} capital activity is thus a stronger high-income signal than any
demographic attribute --- a directly policy-relevant asymmetry between
capital and wage income. Within every stratum, income rises along a
work/status composite (PC1: hours worked 0.56, sex 0.47, education 0.39,
age 0.36). The comparison views quantify the gradients (Task~2): by
occupation, from 3.7\% (\textit{Other-service}) to 47.6\%
(\textit{Exec-managerial}); by education, monotonically from 15.9\%
(high-school) through 41.3\% (Bachelors) to 74.0\% (professional school);
by age, a life-cycle hump peaking at 40.2\% for ages 50--54; by sex, 30.4\%
of men versus 10.9\% of women --- visible in the projections as the
male-coded region of PC1. The QA view (Task~3) tells the analyst to treat
imputed employment rows as low-confidence: occupation imputation carries
mean modal probability 0.297, and imputed rows over-represent majority
classes.

\paragraph{What we learned about large-scale visualization.}
Three lessons stand out. First, at $\sim$50K items every idiom decision
becomes a scalability decision: raw scatterplots saturate, and opacity,
aggregation or projection --- not a bigger canvas --- are what restore
readability. Second, derived data needs its own visualization: both the
DR coordinates and the imputed values are model outputs, and the coursework's
most honest views (\texttt{variables\_used} tooltips, the QA sheet, the
entropy columns) exist to keep those derivations auditable rather than
authoritative; visual analytics that hides its derivations is
indistinguishable from decoration. Third, the perceptual rules are
checkable: palette safety was \emph{computed} (CVD simulation distances),
convergence was gated on $\hat{R}$, and each pipeline stage printed a
verification before the next ran --- treating design guidance as testable
constraints rather than taste made the workbook right on first open.

\begin{thebibliography}{9}
\bibitem{munzner} T.~Munzner, \emph{Visualization Analysis and Design}.
CRC Press, 2014.
\bibitem{sedlmair} M.~Sedlmair, T.~Munzner and M.~Tory,
``Empirical guidance on scatterplot and dimension reduction technique
choices,'' \emph{IEEE Trans.\ Visualization and Computer Graphics},
19(12):2634--2643, 2013.
\bibitem{shneiderman} B.~Shneiderman, ``The eyes have it: a task by data
type taxonomy for information visualizations,'' \emph{Proc.\ IEEE Symp.\
Visual Languages}, 1996.
\bibitem{vandermaaten} L.~van der Maaten and G.~Hinton, ``Visualizing data
using t-SNE,'' \emph{J.\ Machine Learning Research}, 9:2579--2605, 2008.
\bibitem{kohavi} R.~Kohavi, ``Scaling up the accuracy of naive-Bayes
classifiers: a decision-tree hybrid,'' \emph{Proc.\ KDD}, 1996.
\bibitem{cleveland} W.~S.~Cleveland and R.~McGill, ``Graphical perception:
theory, experimentation, and application to the development of graphical
methods,'' \emph{J.\ American Statistical Association}, 79(387):531--554,
1984.
\end{thebibliography}

\end{document}
